\documentclass[10pt,a4paper]{article}
\usepackage{amsmath,amssymb,amsthm,graphicx}
\usepackage{amsmath,amssymb,amsthm,hyperref}
\newtheorem{nttheorem}{Theorem}
\newtheorem{ntlemma}[nttheorem]{Lemma}
\newtheorem{ntdefinition}[nttheorem]{Definition}
\newtheorem{ntclaim}[nttheorem]{Claim}
\newcommand{\sz}{\operatorname{size}_{\rm ABP}}
\newcommand{\Res}{\operatorname{Res}}
\title{Polynomial-size arithmetic branching programs for factors via nonmonic Hensel lifting and reconstruction}
\author{Amit Sinhababu, Thomas Thierauf}
\date{doi:10.4230/LIPIcs.CCC.2020.33; CC BY3.0}
\begin{document}
\maketitle
\noindent\textit{Editorial note.} AmitSinhababu and ThomasThierauf. CCC2020 article33,DOI10.4230/LIPIcs.CCC.2020.33. CC BY3.0.\par\medskip
\noindent\textit{Editorial note.} Faithful complete editable transcription of official CCC2020 CC BY3.0 article33,19-page version. One characteristic-zero ABP factor-closure mechanism, original Theorem10. Complete Section2 ABP/degree/size/resultant/PIT context, complete Section3 coefficient/homogeneous extraction, power-root truncation, multiplicity reduction, random monic transformation, coprime starting-point/resultant shifts, bivariate irreducibility and polynomial linear-system proofs, and complete Section4 nonmonic Hensel12 existence/uniqueness, iteration13/Claims14--15, resultant reconstruction16, explicit coefficient linear system and4.4 size analysis retained. All21 official bibliography entries retained. Field characteristic0, total/x/y degree conditions, edge-based ABP size/substitution bounds, monic original factors versus nonmonic lifts, irreducible/coprime/multiplicity premises and all random PIT/shift probability conditions preserved. Original lifting/degree thresholds/printed indices and coefficients retained without mathematical correction; not an independent proof certification. Main result is the stated size closure; no unqualified uniform factoring algorithm for arbitrary field representation added. Original constructive random preprocessing and univariate-factorization steps remain as authors state them. Standard external determinant/resultant/PIT facts retained at author citation boundary, with all current-paper proof steps included. Journal2021 official47-page CC BY4.0 version cached but not mixed into CCC proof;2023 author update not used. PIT reduction/hardness-vs-randomness direct applications not counted as further tasks. No required figures. Root handles compilation/database/publication.\par\medskip
\noindent\textit{Allowed standard prerequisites (editorial).}\begin{itemize}
\item Resultants and polynomial identity testing. Original Lemma2[20] and exact Theorem3 probability d/|S| retained.
\item ABPs for determinants and univariate Euclidean operations. Original Section2 size/substitution argument and determinant[15] premise retained; full current-author linear-system proof included.
\end{itemize}\medskip
\setcounter{nttheorem}{9}
\begin{nttheorem}[Author Theorem 10]
Let $p$ be a polynomial over a field $\mathbb F$ with characteristic0. For all factors $q$ of $p$, we have
\[\sz(q)\leq\operatorname{poly}(\sz(p)).\]
\end{nttheorem}

\section*{2 Preliminaries}
We consider multivariate polynomials over a field $\mathbb F$ of characteristic0. A polynomial $p$ is called square-free, if for any non-constant irreducible factor $q$, the polynomial $q^2$ is not a factor of $p$.
By $\deg(p)$ we denote the total degree of $p$. Let $x$ and $\mathbf z=(z_1,\ldots,z_n)$ be variables and $p(x,\mathbf z)$ be a $(n+1)$-variate polynomial. Then we can view $p$ as a univariate polynomial $p=\sum_i a_i(\mathbf z)x^i$ over $K[x]$, where $K=\mathbb F[\mathbf z]$. The $x$-degree of $p$ is denoted by $\deg_x(p)$. It is the highest degree of $x$ in $p$. Polynomial $p$ is called monic in $x$, if the coefficient $a_{d_x}(\mathbf z)$ is a nonzero constant, where $d_x=\deg_x(p)$.
By $\operatorname{poly}(n)$ we denote the class of polynomials in $n\in\mathbb N$. We denote by $I=\langle x\rangle$ the ideal of polynomials generated by $x$ over the ring $\mathbb F[x,\mathbf z]$. The $k$-th power of the ideal $I$ is the ideal $I^k=\langle x^k\rangle$.
\paragraph{Computational models}
An arithmetic circuit is a directed acyclic graph, whose leaf nodes are labeled by the variables $x_1,\ldots,x_n$ and various constants from the underlying field. The other nodes are labeled by sum gates or product gates. A node labeled by a variable or constant computes the same. A node labeled by sum or product compute the sum or product of the polynomials computed by nodes connected by incoming edges. The size of an arithmetic circuit is the total number of its edges.
An arithmetic formula is a special kind of arithmetic circuit. A formula has the structure of a directed acyclic tree. Every node in a formula has out-degree at most one. As we can not reuse computations in a formula, it is considered to be weaker than circuits.
An arithmetic branching program (ABP) is a layered directed acyclic graph with a single source node and a single sink node. An edge of an ABP is labeled by a variable or a constant from the field. The weight corresponding to a path from the source to the sink is the product of the polynomials labeling the edges on the path. The polynomial $f(x_1,\ldots,x_n)$ computed by the ABP is the sum of the weights of the all possible paths from source to sink.
The size of an ABP is the number of its edges. The size of the smallest ABP computing $f$ is denoted by $\sz(f)$. The degree of a polynomial computed by an ABP of size $s$ is at most $\operatorname{poly}(s)$.
\paragraph{Properties of ABPs}
Univariate polynomials have small ABPs. Let $p(x)$ be a univariate polynomial of degree $d$. It can be computed by an ABP of size $2d+1$, actually even by a formula of that size.
For univariate polynomials $p(x),q(x)$, the extended Euclidian algorithm computes the gcd $h=\gcd(p,q)$ and also the B\'ezout-coefficents, polynomials $a,b$ such that $ap+bq=h$, where $\deg(a)<\deg(q)$ and $\deg(b)<\deg(p)$. Let $p$ have the larger degree, $d=\deg(p)\geq\deg(q)$. Then clearly also $\deg(h),\deg(a),\deg(b)\leq d$, and consequently, all these polynomials,$p,q,h,a,b$ have ABP-size at most $2d+1$.
Let $p(\mathbf x),q(\mathbf x)$ be multivariate polynomials in $\mathbf x=(x_1,\ldots,x_n)$. For the ABP-size with respect to addition and multiplication, we have
\begin{enumerate}\item $\sz(p+q)\leq\sz(p)+\sz(q)$,\item $\sz(pq)\leq\sz(p)+\sz(q)$.\end{enumerate}
For the sum of two ABPs $B_p,B_q$ one can put $B_p$ and $B_q$ in parallel by merging the two source nodes of $B_p$ and $B_q$ into one new source node, and similar for the two sink nodes. For the product, one can put $B_p$ and $B_q$ in series by merging the sink of $B_p$ with the source of $B_q$.
Another operation is substitution. Let $p(x_1,\ldots,x_n)$ and $q_1(\mathbf x),\ldots,q_n(\mathbf x)$ be polynomials. Let $\sz(q_i)\leq s$, for $i=1,\ldots,n$. Then we have
\[\sz(p(q_1,\ldots,q_n))\leq s\cdot\sz(p).\]
To get an ABP for $p(q_1(\mathbf x),\ldots,q_n(\mathbf x))$, replace an edge labeled $x_i$ in the ABP $B_p$ for $p$ by the ABP $B_{q_i}$ for $q_i$.
It is known that the determinant of a symbolic matrix of dimension $n$ can be computed by an ABP of size $\operatorname{poly}(n)$~\cite{ref15}. By substitution, the entries of the matrix can itself be polynomials computed by ABPs.
\paragraph{Resultant}
Given two polynomials $p(x,\mathbf y)$ and $q(x,\mathbf y)$ in variables $x$ and $\mathbf y=(y_1,\ldots,y_n)$, consider them as polynomials in $x$ with coefficients in $\mathbb F[\mathbf y]$. The resultant of $p$ and $q$ w.r.t.$x$, denoted by $\Res_x(p,q)$, is the determinant of the Sylvester matrix of $p$ and $q$. For the definition of the Sylvester matrix, see~\cite{ref20}. Note that $\Res_x(p,q)$ is a polynomial in $\mathbb F[\mathbf y]$.
Basic properties of the resultant are that it can be represented as a linear combination of $p$ and $q$, and that it provides information about the gcd of $p$ and $q$.
\setcounter{nttheorem}{1}
\begin{ntlemma}[Author Lemma 2, See {\cite{ref20}}]
Let $p(x,\mathbf y)$ and $q(x,\mathbf y)$ be polynomials of degree $\leq d$ and $h=\gcd(p,q)$.
\begin{enumerate}\item $\deg(\Res_x(p,q))\leq4d^2$,\item $\exists u,v\in\mathbb F[x,\mathbf y]\ up+vq=\Res_x(p,q)$,\item $\Res_x(p,q)=0\Longleftrightarrow\deg_x(h)>0$.\end{enumerate}
\end{ntlemma}
Note that the problem whether $\Res_x(p,q)=0$ is a polynomial identity test (PIT), because $\Res_x(p,q)\in\mathbb F[\mathbf y]$. It can be solved in a randomized way by the DeMillo-Lipton-Schwartz-Zippel Theorem (see~\cite{ref4} and the references therein for more details and history of this theorem).
\begin{nttheorem}[Author Theorem 3: Polynomial Identity Test]
Let $p(\mathbf x)$ be an $n$-variate nonzero polynomial of total degree $d$. Let $S\subseteq\mathbb F$ be a finite set. For $\boldsymbol\alpha\in S^n$ picked independently and uniformly at random,
\[\Pr[p(\boldsymbol\alpha)=0]\leq\frac d{|S|}.\]
\end{nttheorem}

\section*{3 Pre-processing Steps and Algebraic Tool Kit}
Before we start the Hensel lifting process, a polynomial should fulfill certain properties that the input polynomial might not have. In this section, we describe transformations of a polynomial that achieve these properties such that ABPs can compute the transformation and its inverse, and factors of the polynomials are maintained.
We also explain how to compute homogeneous components and how to solve linear systems via ABPs. We show how handle the special case when the given polynomial is just a power of an irreducible polynomial.
\subsection*{3.1 Computing homogeneous components and coefficients of a polynomial}
Let $p(x,\mathbf z)$ be polynomial of degree $d$ in variables $x$ and $\mathbf z=(z_1,\ldots,z_n)$. Let $B_p$ be an ABP of size $s$ that computes a polynomial $p$. Write $p$ as a polynomial in $x$, with coefficients from $\mathbb F[\mathbf z]$,
\[p(x,\mathbf z)=\sum_{i=0}^dp_i(\mathbf z)x^i.\]
We show that all the coefficients $p_i(\mathbf z)$ have ABPs of size $\operatorname{poly}(s,d)$.
The argument is similar to Strassen's homogenization technique for arithmetic circuits, an efficient way to compute all the homogeneous components of a polynomial. The same technique can be used for ABPs (see~\cite[Lemma 5.2 and Remark]{ref17}). Here we sketch the proof idea.
Each node $v$ of $B_p$ we split into $d+1$ nodes $v_0,\ldots,v_d$, such that node $v_i$ computes the degree $i$ part of the polynomial computed by node $v$, for $i=0,1,\ldots,d$. Consider an edge $e$ between node $u$ and $v$ in $B_p$.
\begin{itemize}\item If $e$ is labeled with a constant $c\in\mathbb F$ or a variable $z_i$, then we put an edge between $u_i$ and $v_i$ with label $c$ or $z_i$, respectively.\item If $e$ is labeled with variable $x$, then we put an edge between $u_i$ and $v_{i+1}$ with label1.\end{itemize}
The resulting ABP has one source node and $d+1$ sink nodes. The $i$-th sink node computes $p_i(\mathbf z)$.
For each edge of $B_p$ we get either $d$ or $d+1$ edges in the new ABP. Hence, its size is bounded by $s(d+1)$.
\setcounter{nttheorem}{3}
\begin{ntlemma}[Author Lemma 4: Coefficient extraction]
Let $p(x,\mathbf z)=\sum_{i=0}^dp_i(\mathbf z)x^i$ be a polynomial. Then $\sz(p_i)\leq(d+1)\sz(p)$, for $i=0,1,\ldots,d$.
\end{ntlemma}
The technique can easily be extended to constantly many variables. For two variables, consider $p(x,y,\mathbf z)=\sum_{i,j}p_{i,j}(\mathbf z)x^iy^j$. Then, from an ABP of size $s$ for $p$ we get ABPs for the coefficients $p_{i,j}(\mathbf z)$ of size $s(d+1)^2$ similarly as above.
In homogenization, we want to compute the homogeneous components of $p$. That is, write $p(\mathbf z)=\sum_{i=0}^dp_i(\mathbf z)$, where $\deg(p_i)=i$. From an ABP $B_p$ for $p$ we get ABPs for the $p_i$'s similarly as above: In the definition of the new edges, only for constant label, we put the edge from $u_i$ to $v_i$. In case of any variable label $z_j$, we put the edge from $u_i$ to $v_{i+1}$ with label $z_j$. Then the $i$-th sink node computes $p_i(\mathbf z)$. The size is bounded by $s(d+1)$.
\begin{ntlemma}[Author Lemma 5: Homogenization]
Let $p(\mathbf z)=\sum_{i=0}^dp_i(\mathbf z)$ be a polynomial with $\deg(p_i)=i$, for $i=0,1,\ldots,d$. Then $\sz(p_i)\leq(d+1)\sz(p)$, for $i=0,1,\ldots,d$.
\end{ntlemma}
\subsection*{3.2 Computing $q$ from $p=q^e$}
A special case is when the given polynomial $p(\mathbf x)$ is a power of one irreducible polynomial $q(\mathbf x)$, i.e.,$p=q^e$, for some $e>1$. This case is handled separately. Kaltofen~\cite{ref10} showed how to compute $q$ for circuits, ABPs, and arithmetic formulas. Here, we give a short proof from Dutta~\cite{ref5}.
\begin{ntlemma}[Author Lemma 6]
Let $p=q^e$, for polynomials $p(\mathbf x),q(\mathbf x)$. Then $\sz(q)\leq\operatorname{poly}(\sz(p))$.
\end{ntlemma}
\begin{proof}
We may assume that $p$ is nonzero; otherwise the claim is trivial. We want to apply Newton's binomial theorem to compute $q=p^{1/e}$. For this we need that $p(0,\ldots,0)=1$. If this is not the case, we first transform $p$ as follows.
\begin{enumerate}\item If $p(0,\ldots,0)=0$, let $\boldsymbol\alpha=(\alpha_1,\ldots,\alpha_n)$ be a point where $p(\boldsymbol\alpha)\ne0$. By the PIT-Theorem, a random point $\boldsymbol\alpha$ will work, with high probability. Now we shift the variables and work with the shifted polynomial $\widetilde p(\mathbf x)=p(\mathbf x+\boldsymbol\alpha)$. Still,$\widetilde p(0,\ldots,0)$ might be different from1. In this case, we also apply the next item to $\widetilde f$.\item If $p(0,\ldots,0)=a_0\ne0$, then we work with $\widetilde p(\mathbf x)=p(\mathbf x)/a_0$. Then $\widetilde p(0,\ldots,0)=1$.\end{enumerate}
Note that both transformations are easily reversible. Hence, in the following we simply assume that $p(0,\ldots,0)=1$.
By Newton's binomial theorem, we have
\begin{equation}q=p^{1/e}=(1+(p-1))^{1/e}=\sum_{i=0}^\infty\binom{1/e}i(p-1)^i.\tag{1}\end{equation}
Note that $p^{1/e}$ is a polynomial of degree $d=\deg(q)$. Since $p-1$ is constant free, the terms $(p-1)^j$ in the RHS of(1) have degree $>d$, for $j>d$. Thus(1) turns into a finite sum modulo the ideal $\langle\mathbf x\rangle^{d+1}$,
\begin{equation}q=\sum_{i=0}^d\binom{1/e}i(p-1)^i\pmod{\langle\mathbf x\rangle^{d+1}}.\tag{2}\end{equation}
Let $\sz(p)=s$. For the polynomial $Q=\sum_{i=0}^d\binom{1/e}i(p-1)^i$ from(2), we clearly have $\sz(Q)\leq\operatorname{poly}(s)$. Finally, to get $q=Q\bmod\langle\mathbf x\rangle^{d+1}$, we have to truncate the terms in $Q$ with degree $>d$. This can be done by computing the homogeneous components of $Q$ as described in Lemma 5. We conclude that $\sz(q)\leq\operatorname{poly}(s)$.
\end{proof}
\subsection*{3.3 Reducing the multiplicity of a factor}
In the earlier works on bivariate and multivariate polynomial factoring, typically the problem is reduced to factoring a square-free polynomial. This is convenient at various places in the Hensel lifting process. The technique to reduce to the square-free case is via taking the gcd of the input polynomial and its derivative. However, for getting upper bounds on the ABP-size of the factors, we want to avoid gcd-computations, because no polynomial size upper bound for the gcd of two ABPs is known.
We avoid this problem by observing that we do not need the polynomial to be square-free. As we will see, it suffices to have one irreducible factor with multiplicity one, and another coprime factor.
Let $p(\mathbf x)$ be the given polynomial, for $\mathbf x=(x_1,\ldots,x_n)$. The special case that $p$ is a power of one irreducible polynomial we just handled in Section 3.2. Hence, we may assume that $p$ has at least two irreducible factors. So let $p=q^ep'$, where $q$ is irreducible and coprime to $p'$.
Consider the derivative of $p$ w.r.t. some variable, say $x_1$.
\begin{equation}\frac{\partial p}{\partial x_1}=q^{e-1}\left((e-1)\frac{\partial q}{\partial x_1}p'+q\frac{\partial p'}{\partial x_1}\right).\tag{3}\end{equation}
Note that $q$ does not divide the factor $(e-1)\frac{\partial q}{\partial x_1}p'+q\frac{\partial p'}{\partial x_1}$ in(3). Hence, the multiplicity of factor $q$ in $\frac{\partial p}{\partial x_1}$ is reduced by one compared to $p$.
For the ABP-size, we write $p$ as a polynomial in $x_1$, i.e.$p(\mathbf x)=\sum_{i=0}^da_ix_1^i$, where the coefficients $a_i$ are polynomials in $x_2,\ldots,x_n$. By Lemma 4, when $p$ has an ABP of size $s$, then the coefficients $a_i$ can be computed by ABPs of size $s'=s(d+1)$. We observe that then the coefficients of the derivative polynomial $\frac{\partial p}{\partial x_1}=\sum_{i=1}^d ia_ix_1^{i-1}$ have ABPs of size $s'+1$.
We repeat taking derivatives $k=e-1$ times and get $\frac{\partial^kp}{\partial x_1^k}$, which has the irreducible factor $q$ with multiplicity one, as desired.
The coefficients of $\frac{\partial^kp}{\partial x_1^k}$ can be computed by ABPs of size $s'+1$. This yields an ABP of size $\operatorname{poly}(s)$ that computes $\frac{\partial^kp}{\partial x_1^k}$.

\subsection*{3.4 Transforming to a monic polynomial}
Given any polynomial $p(\mathbf z)$ in variables $\mathbf z=(z_1,\ldots,z_n)$, there is a standard trick to make it monic in a new variable $x$ by applying a linear transformation on the variables: for $\boldsymbol\alpha=(\alpha_1,\ldots,\alpha_n)\in\mathbb F^n$, let
\[\tau_{\boldsymbol\alpha}:z_i\mapsto\alpha_ix+z_i,\]
for $i=1,\ldots,n$. Let $p_{\boldsymbol\alpha}(x,\mathbf z)$ be the resulting polynomial. Note that $p$ and $p_{\boldsymbol\alpha}$ have the same degree. We show that $p_{\boldsymbol\alpha}(x,\mathbf z)$ is monic in $x$, for a random transformation $\tau_{\boldsymbol\alpha}$.
\begin{ntlemma}[Author Lemma 7: Transformation to monic]
Let $p(\mathbf z)$ be polynomial of total degree $d$. Let $S\subseteq\mathbb F$ be a finite set. For $\boldsymbol\alpha\in S^n$ picked independently and uniformly at random,
\[\Pr[p_{\boldsymbol\alpha}(x,\mathbf z)\text{ is monic in }x]\geq1-\frac d{|S|}.\]
\end{ntlemma}
\begin{proof}
Consider the terms of degree $d$ in $p$. Let $\boldsymbol\beta=(\beta_1,\ldots,\beta_n)$ such that $|\boldsymbol\beta|=\sum_{i=1}^n\beta_i=d$. We denote the term $\mathbf z^{\boldsymbol\beta}=z_1^{\beta_1}\cdots z_n^{\beta_n}$. Then the homogeneous component of degree $d$ in $p$ can be written as $a_d(\mathbf z)=\sum_{|\boldsymbol\beta|=d}c_{\boldsymbol\beta}\mathbf z^{\boldsymbol\beta}$. Note that $a_d$ is a nonzero polynomial.
Now consider the transformed polynomial $p_{\boldsymbol\alpha}$. We have $\deg_x(p_{\boldsymbol\alpha})=d$ and the coefficient of $x^d$ in $p_{\boldsymbol\alpha}$ is $a_d(\boldsymbol\alpha)=\sum_{|\boldsymbol\beta|=d}c_{\boldsymbol\beta}\boldsymbol\alpha^{\boldsymbol\beta}$. When we pick $\boldsymbol\alpha$ at random,$a_d(\boldsymbol\alpha)$ will be a nonzero constant with probability $\geq1-d/|S|$ by the PIT-Theorem, and in this case $p_{\boldsymbol\alpha}(x,\mathbf z)$ is monic in $x$.
\end{proof}
Given an ABP of size $s$ that computes $p(\mathbf z)$, we can construct another ABP of size $3s$ that computes $p_{\boldsymbol\alpha}(x,\mathbf z)$. For the new ABP replace edge labeled by $z_i$ by the ABP computing $\alpha_ix+z_i$. For each old edge, this requires adding two new edges with labels $\alpha_i$ and $x$.
\subsection*{3.5 Handling the starting point of Hensel lifting}
After doing the above pre-processing steps on the given polynomial $p(\mathbf z)$, we call the transformed polynomial $f(x,\mathbf z)$. We can assume that $f$ of degree $d$ can be factorized as $f=gh$, where $g$ and $h$ are coprime and $g$ is irreducible. In the first step of Hensel lifting, we factorize the univariate polynomial $f(x,0,\ldots,0)\equiv f(x,\mathbf z)\pmod{\mathbf z}$. Now, clearly we have the factorization $f(x,0,\ldots,0)=g(x,0,\ldots,0)h(x,0,\ldots,0)$, but these two factors might not be coprime. In this case we do another transformation.
\paragraph{Author Remark}
Although it would suffice for our purpose to start with two coprime factors, the transformation below produces one irreducible factor.
Let $g_0$ be an irreducible factor of $g(x,0,\ldots,0)$. Then we have for some univariate polynomial $h'_0(x)$ and for $h_0(x)=h'_0(x)h(x,0,\ldots,0)$,
\[g\equiv g_0h'_0\pmod{\mathbf z},\qquad f\equiv g_0h_0\pmod{\mathbf z}.\]
We want that $g_0$ is coprime to $h'_0$ and $h_0$. Directly, this might not be the case because all factors of $f(x,0,\ldots,0)$ might have multiplicity $>1$. However, we argue how to ensure this after a random shift $\boldsymbol\alpha$ of $f$. That is, we consider the function $f(x,\mathbf z+\boldsymbol\alpha)$
\begin{enumerate}
\item First, we show how to achieve that $g_0$ is coprime to $h'_0$.
Since $g$ is irreducible, it is also square-free, and hence,$\gcd(g,\partial g/\partial x)=1$. By Lemma 2, the resultant $r(\mathbf z)=\Res_x(g,\partial g/\partial x)$ is a polynomial of degree $\leq4d^2$ and $r(\mathbf z)\ne0$. Hence, at a random point $\boldsymbol\alpha\in[8d^2]^n$, we have $r(\boldsymbol\alpha)\ne0$ with high probability. At such a point $\boldsymbol\alpha$, we have that $g(x,\boldsymbol\alpha)$ is square-free. Therefore,$g(x,\mathbf z)$ is square-free modulo $(\mathbf z-\boldsymbol\alpha)$, or, equivalently,$g(x,\mathbf z+\boldsymbol\alpha)$ is square-free modulo $\mathbf z$. Hence, when we define $g_0$ and $h'_0$ from $g(x,\mathbf z+\boldsymbol\alpha)$ instead of $g(x,\mathbf z)$, they will be coprime.
\item Similarly, we can achieve that $g_0$ is coprime to $h_0$. By the first item, it now suffices to get $g_0$ coprime to $h(x,0,\ldots,0)$.
For showing this, we use that $g_0$ is coprime to $h'_0$ and prove that $g(x,0,\ldots,0)$ is coprime to $h(x,0,\ldots,0)$. Consider the resultant of $g$ and $h$ w.r.t.$x$, the polynomial $r'(\mathbf z)=\Res_x(g,h)$ has degree $\leq4d^2$. Since $g$ and $h$ are coprime,$r'(\mathbf z)\ne0$. Hence, at a random point $\boldsymbol\alpha\in[8d^2]^n$, we have $r'(\boldsymbol\alpha)\ne0$ with high probability, and hence $g(x,\boldsymbol\alpha)$ and $h(x,\boldsymbol\alpha)$ are coprime univariate polynomials. Therefore,$g(x,\mathbf z)$ and $h(x,\mathbf z)$ are coprime modulo $(\mathbf z-\boldsymbol\alpha)$, or, equivalently,$g(x,\mathbf z+\boldsymbol\alpha)$ and $h(x,\mathbf z+\boldsymbol\alpha)$ are coprime modulo $\mathbf z$.
\end{enumerate}
Combining the two items, a random point $\boldsymbol\alpha\in[8d^2]^n$ will fulfill both properties with high probability. So instead of factoring $f(x,\mathbf z)$, we do a coordinate transformation $\mathbf z\mapsto\mathbf z+\boldsymbol\alpha$ and factor $f(x,\mathbf z+\boldsymbol\alpha)$ instead. From these factors, we easily get the factors of $f(x,\mathbf z)$ by inverting the transformation. Note also that when $f(x,\mathbf z)$ is monic in $x$, the same holds for $f(x,\mathbf z+\boldsymbol\alpha)$.
In the next section, we do another transformation on the input polynomial. We apply a map on the variables that maps $x$ to $x$ and $z_i$ is mapped to $yz_i$, for a new variable $y$ and $i=1,\ldots,n$. Then we factorize the transformed polynomial modulo $y$. Note that in this case, going modulo $y$ has the same effect of going modulo $\mathbf z$. So we can use the above argument to ensure the starting condition for Hensel lifting is satisfied.
\subsection*{3.6 Reducing multivariate factoring to the bivariate case}
Factoring multivariate polynomials can be reduced to the case of bivariate polynomials (see~\cite{ref12}). Let $x,y$ and $\mathbf z=(z_1,\ldots,z_n)$ be variables and let $f(x,\mathbf z)$ be the given polynomial. With $f\in\mathbb F[x,\mathbf z]$, we associate the polynomial $\widehat f\in\mathbb F[x,y,\mathbf z]$ defined by
\[\widehat f(x,y,\mathbf z)=f(x,yz_1,\ldots,yz_n).\]
The point now is to consider $\widehat f$ as a polynomial in $\mathbb F[\mathbf z][x,y]$, that is, as a bivariate polynomial in $x$ and $y$ with coefficients in $\mathbb F[\mathbf z]$. We list some properties.
\begin{enumerate}\item $f(x,\mathbf z)=\widehat f(x,1,\mathbf z)$,\item $\deg(\widehat f)\leq2\deg(f)$,\item $f\text{ monic in }x\Longrightarrow\widehat f\text{ monic in }x$,\item $f=gh\Longrightarrow\widehat f=\widehat g\widehat h$,\item $\widehat f=g'h'\Longrightarrow f=g'(x,1,\mathbf z)h'(x,1,\mathbf z)$.\end{enumerate}
By property 4, factors of $f$ yield factors of $\widehat f$. The following lemma shows that also the irreducibility of the factors is maintained.
\begin{ntlemma}[Author Lemma 8]
Let $f$ be monic in $x$ and $g$ be a non-trivial irreducible factor of $f$. Then $\widehat g$ is a non-trivial irreducible factor of $\widehat f$.
\end{ntlemma}
\begin{proof}
By property 4 above,$\widehat g$ is a factor of $\widehat f$. We argue that $\widehat g$ is irreducible.
Let $\widehat g=uv$ be a factorization of $\widehat g$. By item 5 above, this yields a factorization of $g$ as $g=u(x,1,\mathbf z)v(x,1,\mathbf z)$. Since $g$ is monic in $x$, the same holds for $\widehat g$, and therefore also for factors $u,v$. Hence, either $u$ or $v$ must be constant, because otherwise they would provide a non-trivial factorization of $g$.
\end{proof}
Thus, to get an ABP for an irreducible factor $g$ of $f$, first we show that there is an ABP for the irreducible factor $\widehat g$. This yields an ABP for $g$ by substituting $g=\widehat g(x,1,\mathbf z)$.
Given an ABP $B_f$ of size $s$ for $f$, we get an ABP $B_{\widehat f}$ for $\widehat f$ by putting an edge labeled $y$ in series with every edge labeled $z_i$ in $B_f$, so that $B_{\widehat f}$ computes $yz_i$ at every place where $B_f$ uses $z_i$. Hence, the size of $B_{\widehat f}$ is at most $2s$.
\subsection*{3.7 Solving a linear system with polynomials as matrix entries}
We show how to solve a linear system $M\mathbf v=0$ for a polynomial matrix $M$ with entries from $\mathbb F[\mathbf z]$ given as ABPs. We are seeking for a nonzero vector $\mathbf v$. Note that such a $\mathbf v$ exists over the ring $\mathbb F[\mathbf z]$ iff it exists over the field $\mathbb F(\mathbf z)$.
Except for minor modifications, this follows from classical linear algebra. Kopparty, Saraf, and Shpilka~\cite[Lemma 2.6]{ref12} have shown the same result for circuits. The proof works as well for ABPs.
\begin{ntlemma}[Author Lemma 9: Solving linear systems, {\cite{ref12}}]
Let $M=(m_{i,j}(\mathbf z))_{i,j}$ be a polynomial matrix of dimension $k\times m$ and variables $\mathbf z=(z_1,\ldots,z_n)$, where the entries are polynomials $m_{i,j}\in\mathbb F[\mathbf z]$ that can be computed by ABPs of size $s$.
Then there is an ABP of size $\operatorname{poly}(k,m,s)$ computing a nonzero vector $\mathbf v\in\mathbb F[\mathbf z]^m$ such that $M\mathbf v=0$ (if it exists).
\end{ntlemma}
\begin{proof}
After swapping rows of $M$, we ensure that the $j\times j$ submatrix $M_j$ that consists of the first $j$ rows and the first $j$ columns has full rank, iteratively for $j=1,2,\ldots$.
For $j=1$ this means to find a nonzero entry in the first column and swap that row with the first row. If the first column is a zero-column, then $\mathbf v=(1\ 0\ \cdots\ 0)^T$ is a solution and we are done. To extend from $j$ to $j+1$, suppose we have ensured that $M_j$ has full rank. Now we search for a row from row $j+1$ on, such that after a swap with row $j+1$, the submatrix $M_{j+1}$ has full rank. This can be tested by Lemma 3. If no such row exists, then the process stops at $j$. If $j=m$ then $M$ has full rank and the zero vector is the only solution. Otherwise, assume the above process stops with $j<m$.
Now Cramer's rule can be used to find the unique solution $\mathbf u=(u_1\ u_2\ \cdots\ u_j)^T$ of the system
\[M_j\mathbf u=(m_{1,j+1}\ m_{2,j+1}\ \cdots\ m_{j,j+1})^T.\]
We have $u_i=\det M_j^i/\det M_j$, where $M_j^i$ is the matrix obtained by replacing the $i$-th column of $M_j$ by the vector $(m_{1,j+1}\ \cdots\ m_{j,j+1})^T$. Now, define
\[\mathbf v=(\det M_j^1\ \det M_j^2\ \cdots\ \det M_j^j\ -\det M_j\ 0\ \cdots\ 0)^T.\]
Then $\mathbf v$ is a solution to the original system. Its entries are determinants of matrices with entries computed by ABPs of size $s$. Hence, all the entries of $\mathbf v$ have ABPs of size $\operatorname{poly}(k,m,s)$.
\end{proof}

\section*{4 Factors of Arithmetic Branching Programs}
In this section, we prove that ABPs are closed under factoring.
\setcounter{nttheorem}{10}
We prove Theorem 10 in the rest of this section. First observe that it suffices to prove the $\operatorname{poly}(s)$ size upper bound for the irreducible factors of $p$. This yields the same bound for all the factors. The case when $p=q^e$ is proved in Section 3.2. So it remains to consider the general case when $p=p_1^{e_1}\cdots p_m^{e_m}$, for $m\geq2$, where $p_1,\ldots,p_m$ are the different irreducible factors of $p$. We want to prove an ABP size upper bound for an irreducible factor, say $p_1$.
We start by several transformations on the input polynomial $p(\mathbf z)$, where $\mathbf z=(z_1,\ldots,z_n)$.
\begin{enumerate}
\item As described in Section 3.3, taking $k=e_1-1$ times the derivative w.r.t. some variable, say $z_1$, we get the polynomial $p'(\mathbf z)=\frac{\partial^kp(\mathbf z)}{\partial z_1^k}$, where the factor $p_1$ has multiplicity1.
\item Next, by Lemma 7, we transform $p'(\mathbf z)$ to a polynomial $p''(x,\mathbf z)$ that is monic in $x$, for a new variable $x$. Thereby also the factors of $p'(\mathbf z)$ are transformed, maintaining their irreducibilty and multiplicity. The degree of $p''$ is twice the degree of $p'$.
\item At this point, we may have to shift the variables $\mathbf z$ as described in Section 3.5 to ensure the properties needed for starting the Hensel lifting. This shift preserves the monicness and the irreducibility of the factors.
\item Finally, the transformation to a bivariate polynomial is explained in Section 3.6. This yields polynomial $p'''(x,y,\mathbf z)$, with new variable $y$ and monic in $x$. We rewrite $p'''$ as a polynomial in $x$ and $y$ with coefficients in the ring $K=\mathbb F[\mathbf z]$ and call the representation $f$. That is,$f(x,y)\in K[x,y]$. By Lemma 8, the tranformation maintains irreducible factors. Note also that by the definition of $p'''$ we have $f(x,0)=p'''(x,0,0,\ldots,0)=f(x,y)\bmod y$, so that $f(x,y)\bmod y$ is univariate.
\end{enumerate}
The main part now is to factor $f(x,y)\in K[x,y]$, say $f=gh$, where $g\in K[x,y]$ is irreducible and coprime to $h\in K[x,y]$, and $f,g,h$ are monic in $x$ and have $x$-degree $\geq1$. Let $d$ be the total degree of $f$ in $x,y$.
From the factor $g$ of $f$, we will recover the factor $f_1$ of $p$ by reversing the above transformations. We show that $g$ can be computed by an ABP of size $\operatorname{poly}(s)$. It follows that the irreducible factor $f_1$ has an ABP of size $\operatorname{poly}(s)$.
The basic strategy is to first factor the univariate polynomial $f\bmod y$, and then apply Hensel lifting to get a factorization of $f\bmod y^t$, for large enough $t$. Finally, from the lifted factors modulo $y^t$, we compute the absolute factors of $f$.
\subsection*{4.1 Hensel lifting}
There are various versions of Hensel lifting in the literature (see for example~\cite{ref18}). In our case, an ABP should be able to perform several iterations of the lifting. Therefore we use the lifting in a way suitable for ABPs. In particular, in contrast to other presentations, we will not maintain the monicness of the lifted factors.
Hensel lifting works over rings $R$ modulo an ideal $I\subseteq R$. In our case,$R=K[x,y]$, where $K=\mathbb F[\mathbf z]$, and $I=\langle y\rangle^t$, for some $t\geq1$.
\begin{ntdefinition}[Author Definition 11: Lifting]
Let $R$ be a ring and $I\subseteq R$ be an ideal. Let $f,g,h,a,b\in R$ such that $f\equiv gh\pmod I$ and $ag+bh\equiv1\pmod I$. Then we call $g',h'\in R$ a lift of $g,h$ with respect to $f$, if
\begin{enumerate}
\item[(i)] $f\equiv g'h'\pmod{I^2}$,
\item[(ii)] $g'\equiv g\pmod I$ and $h'\equiv h\pmod I$, and
\item[(iii)] $\exists a',b'\in R\ a'g'+b'h'\equiv1\pmod{I^2}$.
\end{enumerate}
\end{ntdefinition}
\paragraph{Author Remark}
The three conditions in Definition 11 are the invariants when iterating the lifting.
Note that the last condition is actually redundant. It follows from the assumptions together with the second condition. This can be seen in the proof of Lemma 12 below, where a lift $g',h'$ from $g,h$ is constructed, together with $a',b'$. When we show that condition(iii) holds, we do not use the specific form of $g',h'$ constructed there, it suffices to have condition(ii).
\begin{ntlemma}[Author Lemma 12: Hensel Lifting]
Let $R$ be a ring and $I\subseteq R$ be an ideal. Let $f,g,h,a,b\in R$ such that $f\equiv gh\pmod I$ and $ag+bh\equiv1\pmod I$. Then we have:
\begin{enumerate}
\item (Existence). There exists a lift $g',h'$ of $g,h$ w.r.t.$f$.
\item (Uniqueness). For any other lift $g^*,h^*$ of $g,h$ w.r.t.$f$, there exists a $u\in I$ such that
\[g^*\equiv g'(1+u)\pmod{I^2}\quad\text{and}\quad h^*\equiv h'(1-u)\pmod{I^2}.\]
\end{enumerate}
\end{ntlemma}
\begin{proof}
We first show the existence part. Let
\begin{enumerate}\item $e=f-gh$,\item $g'=g+be$ and $h'=h+ae$,\item $c=ag'+bh'-1$,\item $a'=a(1-c)$ and $b'=b(1-c)$.\end{enumerate}
We verify that $g',h'$ is a lift of $g,h$. Because $f\equiv gh\pmod I$, we have $e=f-gh\equiv0\pmod I$. In other words,$e\in I$. It follows that $g'\equiv g\pmod I$ and $h'\equiv h\pmod I$.
Next we show that $f\equiv g'h'\pmod{I^2}$.
\begin{align*}f-g'h'&=f-(g+be)(h+ae)\\&=f-gh-e(ag+bh)-abe^2\\&\equiv e-e(ag+bh)\pmod{I^2}\\&\equiv e(1-(ag+bh))\pmod{I^2}\\&\equiv0\pmod{I^2}.\end{align*}
In the second line, note that $e^2\in I^2$. The last equality holds because $e\in I$ and $1-(ag+bh)\in I$.
Now, we verify that $a'g'+b'h'\equiv1\pmod{I^2}$. First, observe that
\begin{align*}c&=ag'+bh'-1\\&\equiv ag+bh-1\pmod I\\&\equiv0\pmod I.\end{align*}
Hence,$c\in I$ and we conclude that $a'\equiv a\pmod I$ and $b'\equiv b\pmod I$. Now,
\begin{align*}a'g'+b'h'-1&=a(1-c)g'+b(1-c)h'-1\\&=ag'+bh'-1-c(ag'+bh')\\&=c-c(ag'+bh')\\&=c(1-(ag'+bh'))\\&\equiv0\pmod{I^2}.\end{align*}
The last equality holds because $c\in I$ and $1-(ag'+bh')\equiv-c\equiv0\pmod I$.
For the uniqueness part, let $g^*,h^*$ be another lift of $g,h$. Let $\alpha=g^*-g'$ and $\beta=h^*-h'$. By Definition 11(ii), we have $g'\equiv g\equiv g^*\pmod I$ and $h'\equiv h\equiv h^*\pmod I$, and therefore $\alpha,\beta\in I$.
We first show
\begin{equation}\beta g'+\alpha h'\equiv0\pmod{I^2}.\tag{4}\end{equation}
\begin{align*}\beta g'+\alpha h'&=\beta g'+(g^*-g')h'\\&=\beta g'+g^*h'-g'h'\\&\equiv\beta g'+g^*h'-g^*h^*\pmod{I^2}\\&\equiv\beta g'-\beta g^*\pmod{I^2}\\&\equiv-\alpha\beta\pmod{I^2}\\&\equiv0\pmod{I^2}.\end{align*}
Define $u=a'\alpha-b'\beta$. Because $\alpha,\beta\in I$, also $u\in I$. Then, by(4) and because $a'g'+b'h'\equiv1\pmod{I^2}$, we have
\begin{align*}g'(1+u)&=g'(1+(a'\alpha-b'\beta))\\&=g'+a'g'\alpha-b'g'\beta\\&\equiv g'+a'g'\alpha+b'h'\alpha\pmod{I^2}\\&\equiv g'+\alpha\pmod{I^2}\\&\equiv g^*\pmod{I^2}.\end{align*}
Similarly, we get $h^*\equiv h'(1-u)\pmod{I^2}$.
\end{proof}
For the ABP-size, recall that the size just adds up when doing additions or multiplications. Hence, when $f,g,h,a,b$ have ABPs of size $\leq s$ and we construct ABPs for $g',h',a',b'$ according to steps 1--4 in the above proof, then we get ABPs of size $O(s)$.
\paragraph{Author Remark}
In the monic version of Hensel Lifting there is a division in addition to the4 steps from above. When we assume that $g$ is monic, we can compute polynomials $q$ and $r$ such that $g'-g=qg+r$, where $\deg_x(r)<\deg_x(g)$. Then one can show that $\widehat g=g+r$ and $\widehat h=h'(1+q)$ are a lift of $g,h$ w.r.t.$f$. Moreover, when $g,h$ are monic, so are $\widehat g,\widehat h$. Also the B\'ezout-coefficients $\widehat a,\widehat b$ can be computed. For $\widehat c=a\widehat g+b\widehat h-1$, let $\widehat a=a(1-\widehat c)$ and $\widehat b=b(1-\widehat c)$.
The advantage of the monic version is that the result is really unique. There is no $1+u$ factor between monic lifts. The disadvantage is the extra division which would blow up the ABP-size too much.

\subsection*{4.2 Iterating Hensel lifting}
Let $f=gh$, where $g$ is irreducible and coprime to $h$, and $f,g,h$ are monic in $x$ with $x$-degree $\geq1$.
To start the Hensel lifting procedure, we factor the univariate polynomial $f(x,0)=f\bmod y$ as $f(x,0)=g_0(x)h_0(x)$, where $g_0$ is a divisor of $g\bmod y$, and coprime to $h_0$, and $\deg_x(g_0)\geq1$. Recall that by the pre-processing in Section 3.5, we may assume that there is such a decomposition of $f(x,0)$.
By the Euclidian algorithm, there are polynomials $a_0(x),b_0(x)$ such that $a_0g_0+b_0h_0=1$. Hence, for $I_0=\langle y\rangle$, we have $a_0g_0+b_0h_0\equiv1\pmod{I_0}$ and initiate Hensel lifting with
\[f\equiv g_0h_0\pmod{I_0}.\]
We iteratively apply Hensel lifting to $g_0,h_0$ as described in the proof of Lemma 12. Each time, the ideal gets squared. Let $I_k=I_0^{2^k}$. That is, we get polynomials $g_k,h_k$ such that
\[f\equiv g_kh_k\pmod{I_k},\]
and $g_k,h_k$ is a lift of $g_{k-1},h_{k-1}$ w.r.t.$f$. The following lemma states that $g_k$ divides $g$ modulo $I_k$, for all $k\geq0$. In a sense, the $g_k$'s approximate $g$ modulo increasing powers of $y$.
\begin{ntlemma}[Author Lemma 13]
With the notation from above, for all $k\geq0$ and some polynomial $h'_k$,
\[g\equiv g_kh'_k\pmod{I_k}\quad\text{and}\quad h_k\equiv hh'_k\pmod{I_k}.\]
Moreover,$g_k,h'_k$ is a lift of $g_{k-1},h'_{k-1}$ w.r.t.$g$ and $\deg_x(h'_k)\leq\deg_x(h_k)=2^{O(k)}$, for $k\geq1$.
\end{ntlemma}
\begin{proof}
The proof is by induction on $k\geq0$. For the base case, we have that $g_0$ divides $g$ modulo $I_0$, as explained above. Thus, for some polynomial $h'_0$ that is coprime to $g_0$, we have
\[g\equiv g_0h'_0\pmod{I_0}.\]
Hence, we have $h_0\equiv h'_0h\pmod{I_0}$. Note that $h'_0$ can be just1.
For the inductive step, assume that
\begin{equation}g\equiv g_{k-1}h'_{k-1}\pmod{I_{k-1}}\quad\text{and}\quad h_{k-1}\equiv hh'_{k-1}\pmod{I_{k-1}}.\tag{5}\end{equation}
Let $g'_k,h''_k$ be a lift of $g_{k-1},h'_{k-1}$ w.r.t.$g$, so that in particular
\begin{equation}g'_kh''_k\equiv g\pmod{I_k}.\tag{6}\end{equation}
We claim that then $g'_k,hh''_k$ is a lift of $g_{k-1},hh'_{k-1}$, i.e., of $g_{k-1},h_{k-1}$ by(5), w.r.t.$f$.
\begin{ntclaim}[Author Claim 14]
$g'_k,hh''_k$ is a lift of $g_{k-1},h_{k-1}$ w.r.t.$f$.
\end{ntclaim}
\begin{proof}
We check the three conditions for a lift in Definition 11. For the product condition(i), we have by(6)
\[g'_khh''_k=(g'_kh''_k)h\equiv gh\pmod{I_k}.\]
For condition(ii), we have $g'_k\equiv g_{k-1}\pmod{I_{k-1}}$ by assumption and similarly
\[hh''_k\equiv hh'_{k-1}\equiv h_{k-1}\pmod{I_{k-1}}.\]
By the remark after Definition 11, the condition(iii) already follows now. This proves Claim 14.
\end{proof}
Recall that also $g_k,h_k$ is a lift of $g_{k-1},h_{k-1}$. Hence, by the uniqueness property of Hensel lifting, there is a $u\in I_{k-1}$ such that
\begin{equation}g'_k\equiv g_k(1+u)\pmod{I_k}\quad\text{and}\quad hh''_k\equiv h_k(1-u)\pmod{I_k}.\tag{7}\end{equation}
Now observe that we can move the factor $1+u$: we have that $g_k(1+u),hh''_k$ is a lift of $g_{k-1},h_{k-1}$, then also $g_k,hh''_k(1+u)$ is a lift of $g_{k-1},h_{k-1}$.
\begin{ntclaim}[Author Claim 15]
$g_k,hh''_k(1+u)$ is a lift of $g_{k-1},h_{k-1}$ w.r.t.$f$.
\end{ntclaim}
\begin{proof}
We check the conditions for a lift in Definition 11. The first two of them are trivial: moving the factor $1+u$ clearly does not change the product. Because $u\in I_{k-1}$ we still have the equality with the factors $g_{k-1}$ and $h_{k-1}$ modulo $I_{k-1}$, respectively.
By the remark after Definition 11, the third condition already follows, but it also easy to check now:
Let $a,b\in R$ such that $ag_k+bh_k\equiv1\pmod{I_k}$. It follows by(7) that
\[ag_k+bhh''_k(1+u)\equiv ag_k+bh_k(1-u)(1+u)\equiv ag_k+bh_k(1-u^2)\equiv1\pmod{I_k}.\]
This proves Claim 15.
\end{proof}
Now, define $h'_k=h''_k(1+u)$. Note that
\begin{equation}h'_k\equiv h''_k\equiv h'_{k-1}\pmod{I_{k-1}}.\tag{8}\end{equation}
By(7), we have
\begin{equation}hh'_k\equiv hh''_k(1+u)\equiv h_k(1-u)(1+u)\equiv h_k\pmod{I_k}.\tag{9}\end{equation}
By(9) we have
\begin{equation}f=gh\equiv g_kh_k\equiv g_khh'_k\pmod{I_k}.\tag{10}\end{equation}
It follows from(10) that $gh\equiv g_kh'_kh\pmod{I_k}$. Now we want to cancel $h$ in the last equation and conclude that $g\equiv g_kh'_k\pmod{I_k}$. This we can do because $h$ is monic in $x$, it does not contain a factor $y$, i.e.$h\notin I_0$. Hence, together with(8), we conclude that $g_k,h'_k$ is a lift of $g_{k-1},h'_{k-1}$ w.r.t.$g$.
For the $x$-degree of $h'_k$, consider the equation $h_k\equiv hh'_k\pmod{I_k}$. Since $h$ is monic in $x$ the highest $x$-degree term in the product $hh'_k$ will survive the modulo operation. Therefore $\deg_x(h'_k)\leq\deg_x(h)+\deg_x(h'_k)=\deg_x(h_k)$.
To bound the degree of $h_k$ observe that in each round of the Hensel lifting, the maximum possible degree is bounded by a constant multiple($\leq5$) of the maximum degree in the previous round. After $k$ iterations, the degree of the lifted factors is therefore bounded by $2^{O(k)}$.
\end{proof}
For the ABP-size, we observed at the end of Section 4.1 that the size increases by a constant factor in one iteration. Hence, after $k$ iterations, the size increases by a factor $2^{O(k)}$.

\subsection*{4.3 Factor reconstruction for ABP}
We show how to get the absolute factor $g$ of $f$ from the lifted factor. This is called the jump step in Sudan's lecture notes~\cite{ref18}. The difference to the earlier presentations is that our lifted factor might not be monic.
Let $f=gh$, where $f$ has degree $d$, factor $g$ is irreducible and coprime to $h$, and $f,g,h$ are monic in $x$. In the previous section, we started with a factorization $f\equiv g_0h_0\pmod{I_0}$, where $g_0$ is irreducible and coprime to $h_0$. Moreover,$g\equiv g_0h'_0\pmod{I_0}$, for some $h'_0$ such that $h_0=hh'_0\pmod{I_0}$.
Then we apply Hensel lifting, say $t$-times, for some $t$ to be determined below. By Lemma 13, we get a factorization $f\equiv g_th_t\pmod{I_t}$ such that
\begin{equation}g\equiv g_th'_t\pmod{I_t},\tag{11}\end{equation}
for some $h'_t$ such that $h_t\equiv hh'_t\pmod{I_t}$.
Equation(11) gives us a relation between the known $g_t$ and the unknown $g$, via the unknown $h'_t$. We set up a linear system of equations to find a polynomial $\widetilde g\in K[x,y]$ that is monic in $x$ and has minimal degree in $x$ such that
\begin{equation}\widetilde g\equiv g_t\widetilde h\pmod{I_t},\tag{12}\end{equation}
for some polynomial $\widetilde h$. We give some more details to the linear system below, after the next lemma. Before, we show that $\widetilde g$ is indeed the factor we were looking for, for large enough $t$.
\begin{ntlemma}[Author Lemma 16]
$\widetilde g=g$, for $t\geq\log(4d^2+1)$.
\end{ntlemma}
\begin{proof}
Consider the resultant $r(y)=\Res_x(g,\widetilde g)$. We show that $r(y)=0$. Then it follows from Lemma 2 that $g$ and $\widetilde g$ share a common factor with positive $x$-degree. Since $g$ is irreducible it must be a divisor of $\widetilde g$. Since both of them are monic and have the same $x$-degree, we get equality $\widetilde g=g$, up to constant multiples.
To argue that $r(y)=0$, recall from Lemma 2 that the resultant can be written as $r(y)=ug+v\widetilde g$, for some polynomials $u$ and $v$. By(11) and(12), we have
\[ug+v\widetilde g\equiv g_t(uh'_t+v\widetilde h)\pmod{I_t}.\]
Consider $g_t$ and $w=uh'_t+v\widetilde h$ as polynomials in $y$ with coefficients in $x$. Suppose $g_t=c_0(x)+c_1(x)y+\cdots+c_{d'}(x)y^{d'}$. By the properties of Hensel lifting, we have $g_t\equiv g_0\pmod{I_0}$, and therefore $c_0(x)=g_0(x)$. Recall that $g_0$ is non-constant,$\deg(g_0)\geq1$.
Let $j\geq0$ be the least power of $y$ that appears in $w$ and let its coefficient be $w_j(x)$. Suppose for the sake of contradiction that $j<2^t$. Then the least power of $y$ in $g_tw$ is also $j$, and its coefficient is $g_0(x)w_j(x)$, which is a nonzero polynomial in $x$.
The monomials present in $g_0(x)w_j(x)y^j$ cannot be canceled by other monomials in $g_tw$ because they have larger $y$-degree. It follows that $g_tw\bmod I_t$ is not free of $x$. On the other hand,$r(y)\equiv g_tw\pmod{I_t}$ and $r(y)\in K[y]$ is a polynomial with no variable $x$. This is a contradiction.
We conclude that $j\geq2^t$, which means that $w\equiv0\pmod{I_t}$. Hence, we get $r(y)\equiv0\pmod{I_t}$. Recall that $\deg(r)\leq4d^2$. Now we choose $t\geq\log(4d^2+1)$. Then we can conclude that indeed $r(y)=0$.
\end{proof}
\paragraph{Details for setting up the linear system}
Equation(12) can be used to set up a homogeneous system of linear equations. For the degree bounds of the polynomials, let $d_x=\deg_x(g)$ and $d_y=\deg_y(g)$. We may assume that we know $d_x$ and $d_y$. Let $D_x=\deg_x(g_t)$ and $D_y=4d^2$, where $d=\deg(f)$. Let $D'_x=\deg_x(h_t)$. Recall from Lemma 13 that $\deg_x(\widetilde h)\leq D'_x$. Let
\begin{align*}g_t&=\sum_{i\leq D_x,j\leq D_y}c_{i,j}x^iy^j,\\\widetilde g&=x^{d_x}+\sum_{i<d_x,j\leq d_y}r_{i,j}x^iy^j,\\\widetilde h&=\sum_{i\leq D'_x,j\leq D_y}s_{i,j}x^iy^j,\end{align*}
where the coefficients $c_{i,j},r_{i,j},s_{i,j}$ are polynomials in the variables $z_1,\ldots,z_n$. To ensure that $\widetilde g$ is monic, we set the coefficient of $x^{d_x}$ in $\widetilde g$ to be1.
Note that we have an ABP that computes $g_t$. Hence, there are ABPs for computing the coefficients $c_{i,j}$ of $g_t$ by Lemma 4. The coefficients $r_{i,j},s_{i,j}$ of $\widetilde g$ and $\widetilde h$ we treat as unknowns. Equation(12) now becomes
\begin{equation}\sum_{i\leq d_x,j\leq d_y}r_{i,j}x^iy^j\equiv\left(\sum_{i\leq D_x,j\leq D_y}c_{i,j}x^iy^j\right)\left(\sum_{i\leq D'_x,j\leq D_y}s_{i,j}x^iy^j\right)\pmod{y^{2d^2+1}}.\tag{13}\end{equation}
Now we equate the coefficients of the monomials $x^ky^\ell$ on both sides in(13). Then we get $(D_x+D'_x+1)(D_y+1)$ homogeneous linear equations in $d_x(d_y+1)+(D'_x+1)(D_y+1)$ many unknowns $r_{i,j}$ and $s_{i,j}$. This system can be expressed in the form $M\mathbf v=0$, for a matrix $M$ and unknown vector $\mathbf v$.
By Lemma 9, an ABP can efficiently compute a solution vector $\mathbf v$ of polynomials from $\mathbb F[\mathbf z]$. Note that by(11), a nontrivial solution is guaranteed to exist. From $\mathbf v$, we get the coefficients $r_{i,j}$ of $\widetilde g$, and hence of the factor $g$.
\subsection*{4.4 Size Analysis}
We summarize the bound on the ABP-size of the factor computed. Given polynomial $p$ of degree $d_p$ and $\sz(p)=s$. We have seen that the pre-processing transformations yield a polynomial $f$ of degree $d_f\leq2d_p$ and $\sz(f)=\operatorname{poly}(s)$. Then we do $t=\log(2d_f^2+1)$ iterations of Hensel lifting. The initial polynomials $f_0,g_0,h_0$ have ABP-size bound by $2d_f$. Hence, the polynomials after the last iteration have ABP-size bounded by $2^t\operatorname{poly}(s)=\operatorname{poly}(s,d_p)=\operatorname{poly}(s)$.
From the lifted factor we construct the actual factor of $f$. This step involves solving a linear system. We argued that the resulting polynomial $g$ has ABP-size $\operatorname{poly}(s)$.
Finally, we reverse the transformations from the beginning and get a factor of $p$ that has an ABP of size $\operatorname{poly}(s)$. This finishes the proof of Theorem 10.

\begin{thebibliography}{21}
\bibitem{ref1}Peter Bürgisser. The complexity of factors of multivariate polynomials. Foundations of Computational Mathematics, 4(4):369–396, 2004.
\bibitem{ref2}Peter Bürgisser. Completeness and reduction in algebraic complexity theory, volume 7 of Algorithms and Computation in Mathematic. Springer, 2013.
\bibitem{ref3}Chi-Ning Chou, Mrinal Kumar, and Noam Solomon. Closure of VP under taking factors: a short and simple proof. Technical Report arXiv:1903.02366, arXiv, 2019. arXiv:1903.02366.
\bibitem{ref4}Chi-Ning Chou, Mrinal Kumar, and Noam Solomon. Closure results for polynomial factoriza- tion. Theory of Computing, 15(13):1–34, 2019. doi:10.4086/toc.2019.v015a013.
\bibitem{ref5}Pranjal Dutta. Discovering the roots: Unifying and extending results on multivariate poly- nomial factoring in algebraic complexity. Master’s thesis, Chennai Mathematical Institute, 2018.
\bibitem{ref6}Pranjal Dutta, Nitin Saxena, and Amit Sinhababu. Discovering the roots: Uniform closure results for algebraic classes under factoring. In Proceedings of the 50th Annual ACM SIGACT Symposium on Theory of Computing (STOC), pages 1152–1165. ACM, 2018.
\bibitem{ref7}Michael A. Forbes, Amir Shpilka, Iddo Tzameret, and Avi Wigderson. Proof complexity lower bounds from algebraic circuit complexity. In Proceedings of the 31st Conference on Computational Complexity (CCC), page 32. LIPIcs, 2016.
\bibitem{ref8}Valentine Kabanets and Russell Impagliazzo. Derandomizing polynomial identity tests means proving circuit lower bounds. In Proceedings of the 35th ACM Symposium on Theory of Computing (STOC), pages 355–364. ACM, 2003.
\bibitem{ref9}Erich Kaltofen. Uniform closure properties of p-computable functions. In Proceedings of the 18th Annual ACM Symposium on Theory of Computing (STOC), pages 330–337, 1986.
\bibitem{ref10}Erich Kaltofen. Single-factor Hensel lifting and its application to the straight-line complexity of certain polynomials. In Proceedings of the 19th annual ACM Symposium on Theory of Computing (STOC), pages 443–452. ACM, 1987.
\bibitem{ref11}Erich Kaltofen. Factorization of polynomials given by straight-line programs. Randomness and Computation, 5:375–412, 1989.
\bibitem{ref12}Swastik Kopparty, Shubhangi Saraf, and Amir Shpilka. Equivalence of polynomial identity testing and polynomial factorization. computational complexity, 24(2):295–331, 2015.
\bibitem{ref13}Mrinal Kumar and Ramprasad Saptharishi. Hardness-randomness tradeoffs for algebraic computation. Bulletin of EATCS, 3(129), 2019.
\bibitem{ref14}Meena Mahajan. Algebraic complexity classes. In Perspectives in Computational Complexity, pages 51–75. Springer, 2014.
\bibitem{ref15}Meena Mahajan and V Vinay. Determinant: Old algorithms, new insights. SIAM Journal on Discrete Mathematics, 12(4):474–490, 1999.
\bibitem{ref16}Rafael Oliveira. Factors of low individual degree polynomials. computational complexity, 2(25):507–561, 2016.
\bibitem{ref17}Ramprasad Saptharishi. A survey of lower bounds in arithmetic circuit complexity. https: //github.com/dasarpmar/lowerbounds-survey/releases, 2016.
\bibitem{ref18}Madhu Sudan. Algebra and computation. http://people.csail.mit.edu/madhu/FT98/ course.html, 1998. Lecture Notes.
\bibitem{ref19}Joachim von zur Gathen. Hensel and Newton methods in valuation rings. Mathematics of Computation, 42(166):637–661, 1984.
\bibitem{ref20}Joachim von zur Gathen and Jürgen Gerhard. Modern Computer Algebra. Cambridge University Press, 2013.
\bibitem{ref21}Hans Zassenhaus. On Hensel factorization, I. Journal of Number Theory, 1(3):291–311, 1969.
\end{thebibliography}
\end{document}
